\section{Related Work}


Without a standard and consistent arithmetic behavior, matrix multiplication accelerators (MMAs) raise challenges in AI research and development. Deepseek developers \cite{deepseek-ai_deepseek-v3_2024} have reported that the training accuracy of their LLM is significantly degraded due to the insufficient accumulation precision of the FP8 matrix multiplication instructions of the Hopper architecture. This To mitigate the issue, they use additional CUDA Cores to perform full FP32-precision matrix accumulation. PyTorch developers \cite{pytorch_developers_pytorch_2025} have reported that the FP16 and BF16 MFMA instructions of the CDNA2 architecture will flush subnormal numbers to zero, and FP16 DNN training may fail to converge as subnormal values frequently occur during gradient calculation in the backward pass. To mitigate the issue, they cast FP16 to BF16 for computation in the backward pass, which increases the dynamic range but sacrifices precision.

In addition, previous studies have reported numerical inconsistency issues across different hardware \cite{jezequel_estimation_2015}, GPU vendors \cite{zahid_testing_2024}, and different generations of MMAs \cite{li_discovery_2024} due to intricate properties of floating-point arithmetic. Reproducibility of DNNs is an important goal for AI researchers \cite{pham_problems_2020,chen_towards_2022,he_defeating_2025}, but hardware-related numerical inconsistency makes it rather difficult \cite{xie_revealing_2025}.

Although many studies have modeled the arithmetic behavior of existing MMAs, none of them are verified to be bit-accurate. Raihan et al. \cite{raihan_modeling_2019} intuitively model the dot product operation of the Tensor Core as a full binary tree of multiplication and addition operations. However, this is incorrect, as disproved by Hickmann et al. through numerical experiments \cite{hickmann_experimental_2019}. Their numerical experiments indicate several characteristics of the Tensor Core on the Volta architecture. Fasi et al. \cite{fasi_numerical_2021} extend the numerical experiments to the Tensor Cores on the Turing and Ampere architectures, and Li et al. \cite{li_fttn_2024} extend to the Tensor Core on the Hopper architecture and the Matrix Cores on the CDNA1 and CDNA2 architectures. Xie et al. \cite{xie_revealing_2025} invent an algorithm to reveal the floating-point summation order in MMAs. However, these studies only cover a part of the arithmetic behavior of the accelerators, incomplete for a bit-accurate reference model. Moreover, some of their results are incorrect, as reported by Valpey et al.\cite{valpey_smt_2025}. Without bit-level verification, the results of these studies lack reliability.

